Abbreviate
\[
M_n=2^{dJ_z},\qquad K=K_{J_y,J_z,n},\qquad
A=A_{J_y,J_z,n},\qquad V=V_{L,n},\qquad S=S_n=A+\sqrt{nV},
\]
and put \(b_{y,n}=2^{-(s_y+1)J_y}\) and \(b_{z,n}=2^{-s_zJ_z}\).  All constants below depend only on the fixed density bounds, H\"older radius, and wavelet basis.  In particular, they do not depend on the member of \(\mathcal M_{n,d}(s_y,s_z,\rho)\).

First apply the uniform second-order expansion.  Its linear standard-deviation conclusion gives \(\sqrt{nV}\asymp\rho_n\sqrt n\), while the spectral proposition gives \(A\asymp M_n^{1/2}\).  The assumed bias relation is
\[
b_{y,n}^2+b_{z,n}^2=O(S/n).
\]
After subtracting the two displayed biases and multiplying the expansion by \(n\), every term in its remainder scale is therefore \(O(S)\).  Hence
\[
n\{\widehat\theta^{\mathrm{an}}_{J_y,J_z,n}-\theta_n-B^y_{J_y,n}-B^z_{J_z,n}\}
=n\mathbb L_{J_y,J_z,n}+n\mathbb U_{J_y,J_z,n}+o_P(S). \tag{1}
\]

We next correct the high-frequency part of the linear projection.  Let \(s_{a,n}\) denote the exact centered first derivative of the conditional transport functional in arm \(a\), and let \(s^{(2)}_{a,J_y,J_z,n}\) denote the equality-Hessian derivative obtained by applying the projected inverse-elliptic form to \(\Pi_{J_y,J_z}\Delta_n\).  Uniform positivity of \(q_{a,n}\) makes the one-dimensional inverse distribution maps uniformly Lipschitz.  Expanding their defining identities once around \(\bar q_n\), the first omitted Taylor term is a product of the inverse-elliptic size \(\rho_n\) and the density perturbation size \(\|\Delta_n\|_\infty\).  Thus the nonlinear part is bounded in \(L^2(P_{a,n})\) by \(C\rho_n\|\Delta_n\|_\infty\).

For completeness, the projection part follows directly from the coefficient calculation underlying the inverse-elliptic comparison in the spectral proposition.  If \(d_{j_y,k_y,\gamma}\) are the tensor-wavelet coefficients of \(\Delta_n\), the order-minus-two outcome multiplier and the anisotropic H\"older coefficient bounds give
\[
\sum_{j_y>J_y,k_y,\gamma}2^{-2j_y}d_{j_y,k_y,\gamma}^2
\lesssim 2^{-2(s_y+1)J_y}=b_{y,n}^2,
\]
and
\[
\sum_{j_z(\gamma)>J_z,j_y,k_y}2^{-2j_y}d_{j_y,k_y,\gamma}^2
\lesssim 2^{-2s_zJ_z}=b_{z,n}^2.
\]
The fixed upper and lower density bounds make the corresponding Lebesgue, \(R_n\), and arm-law norms equivalent.  Taking square roots and combining the nonlinear and projection pieces proves the corrected bound
\[
\max_{a\in\mathcal A}
\|s_{a,n}-s^{(2)}_{a,J_y,J_z,n}\|_{L^2(P_{a,n})}
\le C\{\rho_n\|\Delta_n\|_\infty+b_{y,n}+b_{z,n}\}. \tag{2}
\]
The two frequency-tail terms in (2) are not multiplied by \(\rho_n\).

For fold \(e\), write \(\bar\Phi_{a,e}=m^{-1}\sum_{i\in I_{a,e}}\Phi_{J_y,J_z}(W_{a,i,n})\), abbreviate \(\Gamma=\Gamma_{J_y,J_z,n}\) and \(\delta=\delta_{J_y,J_z,n}\), and set
\[
Z_{e,n}=\Gamma^{-1/2}\sqrt m\{\bar\Phi_{1,e}-\bar\Phi_{0,e}-\delta\},
\qquad
G_n=2\sum_{e=1}^2\sqrt m\,\delta^\top B_{J_y,J_z,n}\Gamma^{1/2}Z_{e,n}. \tag{3}
\]
Independent randomized-arm sampling and the definition of \(\Gamma\) imply
\[
E Z_{e,n}=0,\qquad E Z_{e,n}Z_{e,n}^\top=I,
\]
and the two vectors are independent across \(e\).  Abbreviate \(v_k=v_{k,J_y,J_z,n}\), \(\nu_k=\nu_{k,J_y,J_z,n}\), and put \(Z_{e,k}=v_k^\top Z_{e,n}\).  With \(\eta_{k,n}=v_k^\top\Gamma^{-1/2}\sqrt m\,\delta\), diagonalization of \(K=\Gamma^{1/2}B\Gamma^{1/2}\) gives
\[
G_n=2\sum_{e=1}^2\sum_k\nu_k\eta_{k,n}Z_{e,k},
\qquad
\operatorname{Var}(G_n)=8\sum_k\nu_k^2\eta_{k,n}^2. \tag{4}
\]
The empirical average of a centered score error with \(L^2\)-norm \(r_n\) has, after multiplication by \(n\), \(L^2\)-norm at most \(C\sqrt n\,r_n\), by independence.  Consequently (2) yields
\[
\|n\mathbb L_{J_y,J_z,n}-G_n\|_2
\le C\sqrt n\{\rho_n\|\Delta_n\|_\infty+b_{y,n}+b_{z,n}\}. \tag{5}
\]
The first term on the right is \(o(S)\), because \(\|\Delta_n\|_\infty\to0\), \(\rho_n\sqrt n\asymp\sqrt{nV}\), and \(S\ge\sqrt{nV}\).  The bias premise and Cauchy--Schwarz give
\[
\sqrt n(b_{y,n}+b_{z,n})
\le\{2n(b_{y,n}^2+b_{z,n}^2)\}^{1/2}=O(\sqrt S). \tag{6}
\]
The resolution conditions inherited from the expansion have \(J_z\to\infty\).  The spectral proposition then gives \(A\asymp2^{dJ_z/2}\to\infty\), hence \(S\to\infty\) and \(O(\sqrt S)=o(S)\).  Equations (5)--(6) therefore prove
\[
\|n\mathbb L_{J_y,J_z,n}-G_n\|_2=o(S),
\qquad
8\sum_k\nu_k^2\eta_{k,n}^2=nV+o(S^2). \tag{7}
\]
This is the required control of the previously omitted high-frequency linear contribution.

We now analyze the quadratic projection.  The spectral proposition, \(M_n\le D_{J_y,J_z}=o(n^2)\), and \(m=n/2\) imply
\[
M_n\to\infty,\qquad \frac{M_n}{m^2}\to0,
\qquad
\frac{\max_k|\nu_k|}{(\sum_k\nu_k^2)^{1/2}}\lesssim M_n^{-1/2}\to0. \tag{8}
\]
Put \(Q_n=n\mathbb U_{J_y,J_z,n}\).  Canonical centering of the diagonal-free within-arm and cross-arm terms writes
\[
Q_n=\sum_{r<s}h_{n,rs}(W_r,W_s), \tag{9}
\]
where the index \(r\) ranges over the independent evaluation observations and each kernel has conditional mean zero in either argument.  We record the contraction bounds needed below.  Compact support of the boundary wavelets leaves only a fixed number of neighboring spatial indices in each product.  Uniform density bounds convert all arm-law integrals to bounded multiples of Lebesgue integrals.  In the outcome coordinate the squared Hessian weights are summable because
\[
\sum_{j_y\ge j_0}2^{j_y}2^{-4j_y}<\infty,
\]
whereas the number of retained covariate coordinates is of order \(M_n\).  Expanding the second and fourth contractions in (9), these facts and \(\operatorname{tr}(K^2)\asymp M_n\) give
\[
\operatorname{Var}(Q_n)=A^2\{1+o(1)\},qquad
\max_r\sum_{s\ne r}E h_{n,rs}^2\lesssim \frac{M_n}{m},qquad
\sum_{r<s}E h_{n,rs}^4\lesssim\frac{M_n^3}{m^2}. \tag{10}
\]
The only fourth-order cyclic contraction not controlled by the last display is bounded by
\[
\frac{\operatorname{tr}(K^4)}{\operatorname{tr}(K^2)^2}
\le\frac{\|K\|_{\mathrm{op}}^2}{\|K\|_{\mathrm{HS}}^2}\lesssim M_n^{-1}. \tag{11}
\]
Thus (10)--(11) vanish after normalization by \(S^4\), since \(S^2\ge A^2\asymp M_n\) and \(M_n/m^2\to0\).

Write \(G_n=\sum_r\ell_{n,r}(W_r)\), order the independent observations, and define
\[
\mathcal D_{n,r}=\ell_{n,r}(W_r)+\sum_{s<r}h_{n,sr}(W_s,W_r).
\]
Canonicality in (9) makes \((\mathcal D_{n,r})_r\) a martingale-difference array whose sum is \(G_n+Q_n\).  The bounded one-dimensional inverse-map score, (8), and (10)--(11) imply, for each \(\epsilon>0\),
\[
\sum_rE\left[\left(\frac{\mathcal D_{n,r}}S\right)^2
\mathbf1\{|\mathcal D_{n,r}|>\epsilon S\}\right]\longrightarrow0. \tag{12}
\]
Indeed, the fourth-moment expansion of the left side after Markov's inequality is bounded by a constant times \(m^{-1}+M_n/m^2+M_n^{-1}\).  Expanding the paired edges once more, with mixed linear--quadratic contractions bounded by Cauchy--Schwarz and (7), gives
\[
\frac{E\left[\left\{\sum_rE(\mathcal D_{n,r}^2\mid\mathcal F_{r-1})
-E(G_n+Q_n)^2\right\}^2\right]}{S^4}
\lesssim m^{-1}+\frac{M_n}{m^2}+M_n^{-1}+o(1)\longrightarrow0. \tag{13}
\]
Orthogonality of the first and second Hoeffding projections, (7), and (10) also show
\[
E(G_n+Q_n)^2=A^2+nV+o(S^2). \tag{14}
\]

Here is a direct characteristic-function conclusion from (12)--(14).  For fixed \(t\), apply
\[
|e^{ix}-1-ix+x^2/2|\le Cx^2\min\{|x|,1\}
\]
conditionally to each martingale increment.  The sum of the remainders tends to zero by (12); replacing the sum of conditional variances by (14) costs \(o(1)\) by (13).  Iteration over \(r\) therefore yields
\[
E\exp\{it(G_n+Q_n)/S\}
-\exp\left[-\frac{t^2}{2S^2}\{A^2+nV+o(S^2)\}\right]\longrightarrow0. \tag{15}
\]

It remains to compare (15) with the declared finite-sample proxy.  Define
\[
\mathcal Q_n^G=\sum_{e=1}^2\sum_k\nu_k
\{(\xi_{e,k}+\eta_{k,n})^2-(1+\eta_{k,n}^2)\}.
\]
For \(c_{e,k}=\nu_k/S\) and \(l_{e,k}=2\nu_k\eta_{k,n}/S\), direct integration against a standard Gaussian density gives
\[
\log E e^{it\mathcal Q_n^G/S}
=\sum_{e,k}\left[-itc_{e,k}-\frac12\log(1-2itc_{e,k})
-\frac{t^2l_{e,k}^2}{2(1-2itc_{e,k})}\right]. \tag{16}
\]
Equations (7)--(8) give
\[
\sum_{e,k}l_{e,k}^2=\frac{nV+o(S^2)}{S^2},qquad
\max_k\frac{|\nu_k|}{S}\le\frac{\max_k|\nu_k|}{A}\longrightarrow0. \tag{17}
\]
Taylor expansion of the two analytic terms in (16), uniformly for fixed \(t\), now gives
\[
\left|\log E e^{it\mathcal Q_n^G/S}
+\frac{t^2(A^2+nV)}{2S^2}\right|
\le C_t\max_k\frac{|\nu_k|}{S}
\left\{\frac{A^2}{S^2}+\sum_{e,k}l_{e,k}^2\right\}+o(1)=o(1). \tag{18}
\]
Thus (15) and (18) have the same Gaussian characteristic function, with variance
\[
\sigma_n^2=\frac{A^2+nV}{S^2}\in[1/2,1]. \tag{19}
\]
Along every subsequence, compactness of \([1/2,1]\) supplies a further subsequence on which \(\sigma_n^2\) converges.  Equations (15) and (18), together with their uniformly bounded second moments, then imply that both laws converge on that further subsequence to the same centered normal law.  The subsequence criterion for the bounded-Lipschitz metric proves
\[
d_{\mathrm{BL}}\{\mathcal L((G_n+Q_n)/S),\mathcal L(\mathcal Q_n^G/S)\}\to0
\]
and
\[
d_{\mathrm{BL}}\{\mathcal L((G_n+Q_n)/S),\mathcal N(0,\sigma_n^2)\}\to0. \tag{20}
\]
Equations (1) and (7) change the normalized statistic by \(o_P(1)\); truncating at a fixed level and using the Lipschitz property of the test functions shows directly that this changes either bounded-Lipschitz distance by \(o(1)\).  Hence (20) proves the first two conclusions of the theorem with the stated \(\nu_k\), \(\eta_{k,n}\), and \(A^2=4\sum_k\nu_k^2\).

At conditional equality, \(\Delta_n=0\), so \(V=0\), every \(\eta_{k,n}=0\), and \(\mathbb L_{J_y,J_z,n}=0\).  Taking \(S=A\) in (16)--(18) gives
\[
\frac{n\mathbb U_{J_y,J_z,n}}{A_{J_y,J_z,n}}\Rightarrow\mathcal N(0,1).
\]
This Gaussianization is forced by (8); it is not a fixed-rank chaos limit.

Suppose next that \(D_{J_y,J_z}\log n/m=o(1)\).  The observable-pilot lemma gives, for each fold,
\[
\|\widehat K^{(-e)}-K\|_{\mathrm{HS}}=o_P(\|K\|_{\mathrm{HS}}).
\]
Conditionally on the pilots, if \(Z\) is standard Gaussian and \(E_e=\widehat K^{(-e)}-K\), direct Gaussian fourth-moment expansion gives
\[
E\{Z^\top E_eZ-\operatorname{tr}(E_e)\}^2=2\|E_e\|_{\mathrm{HS}}^2=o_P(A^2). \tag{21}
\]
The linear part of the proxy is the first Hoeffding projection already represented by \(G_n\), rather than part of the kernel replacement.  Thus replacing the population quadratic kernel by the average observed quadratic kernel changes the normalized proxy by \(o_P(1)\), which proves the stated observed-kernel clause.

Finally consider the separate fixed-unequal-law assertion.  On the compact support, positivity of both densities bounds the exact centered Kantorovich scores.  For their independent triangular sum, the same conditional expansion of \(e^{itx}\) used above has a remainder bounded by the maximum normalized summand, which is \(O(n^{-1/2})\); hence
\[
\frac{\mathbb L_{J_y,J_z,n}}{\sqrt{V_{L,n}/n}}\Rightarrow\mathcal N(0,1). \tag{22}
\]
The assumed relation \(\sqrt{nV_{L,n}}/A_{J_y,J_z,n}\to\infty\) and the quadratic variance formula imply \(\mathbb U_{J_y,J_z,n}=o_P(\sqrt{V_{L,n}/n})\).  The stated bias and expansion-remainder hypotheses remove the deterministic and residual terms at the same scale.  Finally, \(\widehat V_{L,n}/V_{L,n}\to_P1\) changes the denominator in (22) by a multiplicative \(1+o_P(1)\).  This proves the studentized fixed-unequal-law central limit theorem and completes all claims.